\section{Introduction}
\label{sec:introduction}

\subsection{Problem and Motivation}

Live student showcases historically run on a triplet of disjoint tools:
slide-deck submissions in email, a printed scoring sheet for judges, and a
hand-counted paper ballot for the audience. The cost of this fragmentation is
low only on paper. In practice it produces three failure modes that recur at
every event:

\begin{enumerate}[leftmargin=*]
  \item \textbf{Coordination cost.} The MC, the judges, and the audience all
        consume the same schedule, but each from a different artefact, so a
        five-minute slip in the agenda silently desynchronises everyone.
  \item \textbf{Integrity blind spots.} Paper ballots are easy to stuff and
        impossible to audit. Spreadsheet rubrics are easy to overwrite and
        leave no record of who changed what when.
  \item \textbf{Result latency.} Aggregating scores by hand pushes the
        announcement to the very end of the show, and any tiebreaker dispute
        forces a recess.
\end{enumerate}

The Soft Skills 2026 event at the LAU exposed these failure modes again at the
2025 dry run and asked, in essence: can we replace all three with one piece of
software that does \emph{simultaneous} judging and voting, exposes its work to
the admin in real time, and refuses to be tampered with?

\subsection{Goals}

Spotlight is the answer that this report describes. The goals scoped at the
outset were:

\begin{itemize}[leftmargin=*]
  \item \textbf{Single source of truth.} One database row per project, per
        judge grade, per audience vote, per audit event. Everything else is a
        view over those rows.
  \item \textbf{Real-time push.} The audience and the admin should see state
        change within a second of the change being committed; not in the
        next polling tick.
  \item \textbf{Defensible integrity.} Voting should resist the obvious abuse
        vectors (bots, NAT subnets, browser-token reuse) without requiring an
        account for every audience member.
  \item \textbf{Operable in a small budget.} The whole stack should run on
        free-tier hosts (Vercel, Render) and a free-tier Supabase project,
        because that is the budget a student org actually has.
  \item \textbf{Defensible at oral defence.} Every non-functional concern
        the supervisor is likely to raise should map to a concrete artefact
        in the codebase.
\end{itemize}

\subsection{Non-goals}

What Spotlight is \emph{not} is also worth stating, both to bound the work
and to frame the future-work section:

\begin{itemize}[leftmargin=*]
  \item Not a registration / ticketing system. The event team uses an
        independent flow for that.
  \item Not a content management system. Project descriptions and posters are
        seeded by the organisers and edited via the database, not via an
        in-app editor.
  \item Not a payment system. There is no monetary flow; the only ``value''
        is the rank-ordered list of winners.
\end{itemize}

\subsection{Approach in one paragraph}

Spotlight is a React 19 single-page app served as a static bundle from
Vercel's edge CDN. It talks to an Express 4 API hosted on Render. The API
fronts a Supabase Postgres database for durable state and an in-process LRU
cache for hot reads. Real-time updates use Server-Sent Events over a
long-lived HTTP connection; the client maintains a watchdog reconnect and a
60-second polling fallback. The Anthropic API generates an advisory
``quick read'' for each project; the system prompt is marked
\texttt{cache\_control: ephemeral}, exploiting Anthropic's prompt-caching
feature to reduce per-call cost. An anomaly detector scores every vote on
arrival, and an append-only audit log captures every privileged write.

\subsection{Structure of this report}

Section~\ref{sec:background} surveys related work and the libraries we lean
on. Section~\ref{sec:requirements} states the functional and non-functional
requirements as a contract. Section~\ref{sec:architecture} presents the
component and deployment architecture, with explicit decision records.
Section~\ref{sec:datamodel} defines the schema. Section~\ref{sec:design}
zooms into the four flows that carry most of the system: voting, grading,
real-time session sync, and AI insight generation. Section~\ref{sec:impl}
walks through the implementation with annotated code excerpts.
Section~\ref{sec:nonfunc} addresses the six non-functional dimensions in
turn. Sections~\ref{sec:test} through \ref{sec:future} cover testing,
deployment, and what we would do next.
